\documentclass[12pt,a4paper]{article}
% Preambule commun aux comptes rendus CR_*.tex.

\usepackage{array}
\newcolumntype{C}[1]{>{\centering\arraybackslash}m{#1}}
\usepackage{multirow}
\usepackage{geometry}
\geometry{a4paper, margin=1in}
\usepackage{graphicx}
\usepackage{float}
\usepackage{tikz}
\usetikzlibrary{arrows.meta,positioning}
\usepackage{pgfgantt}
\usepackage{pdflscape}
\usepackage{enumitem}
\usepackage{booktabs}
\usepackage{longtable}
\usepackage{ragged2e}
\usepackage{tabularx}
\usepackage{xparse}
\usepackage{ltablex}
\usepackage{xltabular}
\usepackage{subcaption}

\newcolumntype{Y}{>{\RaggedRight\arraybackslash}X}

\newcommand{\StyledTableSetup}{%
  \footnotesize
  \renewcommand{\arraystretch}{1.2}%
  \setlength{\tabcolsep}{4pt}%
}


% ============================================================
% IHEDN COVER TEMPLATE (XeLaTeX)
% ============================================================
% Compilation (from project root):
%   latexmk -pdf main.tex
%   LATEXMK_SHELL_ESCAPE=0 latexmk -pdf main.tex
%
% Optional environment controls from latexmkrc:
%   LATEXMK_ENGINE=xelatex        (default/enforced)
%   LATEXMK_SHELL_ESCAPE=1|0      (default 1, needed for SVG conversion)
%
% Required package for this template:
%   \usepackage{ihedn-cover}
%
% Note: ihedn-cover already loads needed internal packages
% (fontspec, graphicx, tikz, ragged2e, optional svg).
% ============================================================

% ----------------------------------------------------------------
% Package options (documented)
% ----------------------------------------------------------------
% Geometry scope:
%   apply-geometry            -> force A4 + zero margins while cover is built
%   no-apply-geometry         -> do not alter host geometry (default)
%
% Typesetting freeze:
%   global-typesetting        -> apply freeze globally (intrusive)
%   no-global-typesetting     -> no global freeze (default)
%
% Small-caps handling:
%   local-smallcaps-fix       -> \textsc -> \MakeUppercase only inside cover (default)
%   no-local-smallcaps-fix    -> disable local fix
%   global-smallcaps-fix      -> global override (intrusive)
%   no-global-smallcaps-fix   -> disable global fix (default)
%
% Main font handling:
%   global-mainfont           -> set Times New Roman globally (default)
%   local-mainfont            -> set Times only during cover rendering
%   no-mainfont-override      -> keep host document main font
%
% SVG support:
%   svg-support               -> enable SVG logos (default, needs shell-escape)
%   no-svg-support            -> disable SVG path handling (pdf/png only)
%
% Debug overlay:
%   debug                     -> show mm grid + bounding boxes on cover
%   no-debug                  -> hide debug overlay (default)
% ----------------------------------------------------------------
\usepackage[
  no-apply-geometry,
  no-global-typesetting,
  global-smallcaps-fix,
  local-smallcaps-fix,
  global-mainfont,
  svg-support,
  no-debug
]{ihedn-cover}

% ----------------------------------------------------------------
% Bibliography stack (BibLaTeX/Biber, style from subtree vendor/)
% ----------------------------------------------------------------
\usepackage{polyglossia}
\setdefaultlanguage{french}
\makeatletter
\g@addto@macro\captionsfrench{%
  \renewcommand{\tablename}{Tableau}%
  \renewcommand{\figurename}{Figure}%
}
\makeatother
\usepackage{csquotes}
\usepackage{xurl}
\usepackage{hyperref}
\makeatletter
\let\ihedn@orig@input@path\input@path
\def\input@path{{vendor/biblatex-sciences-po-fr/style/}}
\makeatother
\usepackage[
  backend=biber,
  style=sciencespo,
  hyperref=true,
  autolang=other
]{biblatex}
\AtBeginBibliography{\sloppy\setlength{\emergencystretch}{3em}}
\makeatletter
\let\input@path\ihedn@orig@input@path
\makeatother
\addbibresource{bibliographies/bib/rapport.bib}

% ----------------------------------------------------------------
% tcolorbox
% ----------------------------------------------------------------
\usepackage[most]{tcolorbox}
\tcbuselibrary{breakable}

% Example: use Arial for the full report + cover.
% Uncomment and switch package option to no-mainfont-override.
% \setmainfont{Arial}[
%   UprightFont    = *,
%   BoldFont       = * Bold,
%   ItalicFont     = * Italic,
%   BoldItalicFont = * Bold Italic
% ]

%==================================================
% Liste des propositions
%==================================================

\makeatletter

\newcommand{\listpropositionname}{Liste des propositions}

\newcommand{\listofpropositions}{
  \section*{\listpropositionname}
  \addcontentsline{toc}{section}{\listpropositionname}
  \@starttoc{lop}
}

\newcommand*\l@proposition{\@dottedtocline{1}{1.5em}{2.3em}}

\makeatother

%==================================================
% Compteur des propositions
%==================================================

\newcounter{proposition}

%==================================================
% Style graphique des propositions
%==================================================

\newtcolorbox{propositionbox}[2][]{
  lower separated=false,
  colback=white!80!gray,
  colframe=white!20!black,
  fonttitle=\bfseries,
  colbacktitle=white!30!gray,
  coltitle=black,
  enhanced,
  sharp corners,
  attach boxed title to top left={
    xshift=0.5cm,
    yshift=-2mm
  },
  title=#2,
  #1
}

%==================================================
% Environnement proposition
%==================================================

\newenvironment{proposition}[1]
{
  \refstepcounter{proposition}

  \addcontentsline{lop}{proposition}{%
    \protect\numberline{\theproposition}#1%
  }

  \begin{propositionbox}
  {Proposition~\theproposition~: #1}
}
{
  \end{propositionbox}
}

% Renommage Liste des figures en Table des figures
\addto\captionsfrench{%
  \renewcommand{\listfigurename}{Table des figures}
}

\newcommand{\heure}[2]{{\NoAutoSpacing #1:#2}}

\DeclareFieldFormat[misc]{title}{#1}

% Activer ce package avec l'option none retire la césure.
%\usepackage[none]{hyphenat}

\usepackage{adjustbox}
\usepackage{pdfpages}

\makeatletter
\newcommand{\IhednSetPdfPageCount}[2]{%
  \begingroup
  \edef\ihedn@pdfpath{#2}%
  \ifXeTeX
    \global#1=\XeTeXpdfpagecount "\ihedn@pdfpath" %
  \else
    \ifLuaTeX
      \saveimageresource{\ihedn@pdfpath}%
      \global#1=\lastsavedimageresourcepages\relax
    \else
      \pdfximage{\ihedn@pdfpath}%
      \global#1=\pdflastximagepages\relax
    \fi
  \fi
  \endgroup
}
\makeatother

\newcount\IhednIncludedPdfPageCount
\newcommand{\IhednRenderPdfPagesInBoxes}[1]{%
  \IhednSetPdfPageCount{\IhednIncludedPdfPageCount}{#1}%
  \foreach \IhednIncludedPdfPage in {1,...,\the\IhednIncludedPdfPageCount}{%
    \begin{tcolorbox}[
      colback=gray!5,
      colframe=black,
      boxrule=0.5pt,
      arc=3mm
    ]
    \centering
    \includegraphics[page=\IhednIncludedPdfPage,width=0.95\textwidth]{#1}
    \end{tcolorbox}
    \ifnum\IhednIncludedPdfPage<\IhednIncludedPdfPageCount
      \clearpage
    \fi
  }%
}

\usepackage{tablefootnote}


\begin{document}

% ----------------------------------------------------------------
% Cover keys (all available keys are shown below)
% ----------------------------------------------------------------
% Header block (top-right), logo, title, report lines, subtitle, committee.
\PageDeGardeIHEDN{
    header_1={Union-IHEDN},
    header_2={Association des auditeurs IHEDN},
    header_3={région Paris / Île de France},
    date={le 14 mars 2026},
    logo_path={Figures/logos_IHEDN/Logo_AR_16_PARIS_IDF.jpg},
    title={Crises majeures : quel rôle pour les citoyens engagés et les associations IHEDN ?},
    title_max_lines=2,
    title_fontsize={26.000pt},
    title_baselineskip={28.667pt},
    reportline_1={Rapport du Comité d'étude de l'Association régionale des auditeurs},
    reportline_2={de l'IHEDN région Paris / Île de France (AR 16)},
    subtitle={Compte rendu de la réunion du 05 Mars 2026},
    committee={Comité d'étude, d'octobre 2025 à juin 2026, composé de : Mme Valérie Arlé-Faivre, \textit{présidente} ; M. Aurélien Duvignac-Rosa, \textit{rapporteur} ; M. Pascal Bayot Duponchel ; M. Jean-Jacques Cleynen ; Mme Valérie Cowan ; M. Alain Fontan ; M. Pierre-Marie Giard ; Mme Muriel Joyeux ; M. Bruno Leray ; M. Yves-Henri Renhas ; M. Marc Roure ; Mme Isabelle de Segonzac ; M. Gérard Turck ; Mme Marie Danielle Vasquez-Duchêne.}
}

\section{Éléments de rappel}

La réunion du 5 mars 2026 s'inscrit dans la continuité des réflexions engagées
par le comité d'études autour du rôle potentiel des membres de l'AR~16 dans
le cadre de crises majeures.

Les échanges antérieurs ont mis en évidence que l'engagement des auditeurs
IHEDN ne peut être envisagé qu'en tenant compte de la temporalité propre aux
situations de crise.

À cet égard, trois moments distincts peuvent être identifiés :

\begin{itemize}
    \item \textbf{Avant la crise} : les membres de l'AR~16 peuvent jouer un rôle pertinent en matière de préparation, de sensibilisation et de diffusion d'une culture de résilience et d'esprit de défense au niveau territorial ;
    
    \item \textbf{Pendant la crise} : la gestion opérationnelle étant strictement hiérarchisée et placée sous l'autorité des pouvoirs publics, les auditeurs se mettent prioritairement au service des administrations, institutions ou corps auxquels ils appartiennent ;
    
    \item \textbf{Après la crise} : les membres peuvent contribuer aux retours d'expérience (\textit{RETEX}), à l'analyse sociologique, organisationnelle ou informationnelle des crises, ainsi qu'à la diffusion des enseignements tirés.
\end{itemize}

Dans ce cadre, les échanges ont conduit à formuler une interrogation centrale :
\textbf{quelle est la capacité réelle de mobilisation de l'AR~16 en cas de crise majeure ?}

Cette question renvoie à une problématique plus large concernant la nature
même de l'association : constitue-t-elle avant tout un collectif organisé,
un réseau informel d'expertises, un hub d'agrégation de compétences
individuelles ou encore une porte d'entrée permettant aux autorités
publiques d'identifier certaines ressources issues de la société civile ?

Les discussions ont également souligné une vulnérabilité potentielle :
dans l'hypothèse d'une crise majeure, une partie des membres pourrait être
mobilisée par leurs propres obligations professionnelles ou institutionnelles,
ce qui pourrait limiter la capacité d'action collective de l'association
en tant que telle.

Dans ce contexte, il est apparu nécessaire de disposer d'une connaissance
plus fine de la composition sociologique et professionnelle de l'association.
La réalisation d'une enquête rigoureuse et structurée apparaît comme un
outil particulièrement adapté pour établir un diagnostic crédible.

La mise en place d'un questionnaire, évoquée lors des réunions précédentes,
présente à cet égard un double intérêt :

\begin{itemize}
    \item mieux connaître les profils, compétences et disponibilités des auditeurs ;
    \item raviver le lien entre les membres en manifestant l'attention portée
    par l'association à leur parcours et à leurs expertises.
\end{itemize}

Au-delà de cet objectif de connaissance interne, l'enquête pourrait également
contribuer à clarifier la capacité réelle de mobilisation de l'AR~16 dans le
champ plus large de la résilience nationale.

\newpage

Cette démarche s'inscrit par ailleurs dans un contexte institutionnel marqué
par un intérêt croissant pour la contribution des citoyens engagés et des
réseaux associatifs aux dispositifs de gestion des crises majeures. La lettre
adressée par l'IHEMI le 30 décembre 2025 relative aux enjeux liés aux risques
et aux crises illustre cette dynamique\footcite{IHEMI:2025:LIREC75}.

Enfin, les échanges ont souligné que le cinquantième anniversaire de
l'association constitue une opportunité symbolique forte pour engager une
telle réflexion collective : l'enquête permettrait à la fois de célébrer
l'histoire de l'association et d'interroger sa capacité d'adaptation face aux
défis contemporains.

Dans cette perspective, la réalisation d'un questionnaire pourrait constituer
le point de départ d'un diagnostic partagé sur les ressources, les compétences
et les formes d'engagement potentielles des membres de l'AR~16.

\section{Aval du CODIR}

La mise en œuvre d'une telle démarche suppose toutefois l'aval préalable
du Comité directeur.

À cette fin, Alain Fontan a rédigé une lettre à destination du CODIR afin
de proposer officiellement l'engagement de cette initiative.

Cette lettre, présentée en annexe \ref{annexe:lettre_Alain_Fontan}, a été
adressée le 6 mars 2026. Le Comité directeur, qui se réunira le 17 mars 2026, 
sera amené à échanger sur cette proposition, comme l'a indiqué la présidente 
de l'association, Madame Coralie Noël, dans son mail du 08/03/2026.

Au-delà de la dimension méthodologique, la mise en place de ce questionnaire
conduira nécessairement l'association à s'interroger plus largement sur son
positionnement et sa fonction au sein de l'écosystème de la résilience
territoriale.

Il s'agira notamment de déterminer dans quelle mesure l'AR~16 peut constituer
un acteur de diffusion de l'esprit de défense, un réseau d'expertises
mobilisables ou encore un espace de réflexion collective susceptible
d'apporter une contribution utile à l'analyse et à la compréhension des
crises contemporaines.

\newpage

\appendix

\section{Lettre rédigée par Alain Fontan à destination du Comité Directeur}
\label{annexe:lettre_Alain_Fontan}

La lettre reproduite ci-dessous présente les objectifs généraux de la
démarche ainsi que les conditions de sa mise en œuvre.

\begin{tcolorbox}[
breakable,
colback=gray!5,
colframe=black,
boxrule=0.5pt,
arc=3mm
]
\setlength{\parskip}{0.5em}
IHEDN AR 16

Comité d'études « Crises majeures »

\hfill Paris, le 6 mars 2026

Chère Présidente,\\
Chères et Chers Membres CoDir,

Le 50e anniversaire de notre Association constitue un moment important de son histoire. Au-delà de la célébration légitime du parcours accompli, cet anniversaire peut aussi devenir une occasion de projection vers l'avenir.

Dans un contexte marqué par la multiplication des crises, les réseaux d'expertise et de confiance constituent une ressource précieuse pour la société. L'AR 16 représente, à cet égard, un capital humain et relationnel exceptionnel, encore insuffisamment cartographié.

Notre Comité d'études pose la question : que peut apporter notre réseau d'auditeurs en situation de crise majeure ?

Pour répondre à cette interrogation, nous proposons le lancement, à l'occasion du 50e anniversaire, d'une initiative structurante :

\begin{itemize}
    \item la cartographie des compétences, disponibilités et engagements des membres de l'AR 16.
\end{itemize}

Cette démarche prendra la forme d'une enquête auprès des membres, permettant de :

\begin{itemize}
    \item mieux connaître les profils et expertises des auditeurs ;
    \item identifier la disponibilité et le mode d'engagement en cas de crise ;
    \item structurer des groupes de réflexion ou d'analyse ad hoc.
\end{itemize}

Ainsi, notre Association pourrait contribuer à renforcer la culture de coopération et de résilience, sans jamais se substituer aux responsabilités des acteurs publics.

L'enquête - sous la forme d'un questionnaire anonyme - devient un instrument de connaissance du capital stratégique de l'AR 16. Elle permet aussi de transformer l'expérience individuelle en connaissance collective.

Il ne s'agit ni d'un contrôle, ni d'une obligation, ni d'une évaluation. Les réponses seront traitées de manière anonyme, confidentielle et analysées globalement.

Pour nous accompagner dans cette démarche, le Comité d'études a échangé avec le Dr Victor Violier, sociologue à l'IRSEM. La faisabilité du projet suppose toutefois le concours d'un délégué à la protection des données afin de garantir la conformité au règlement européen sur la protection des données.

Dans cette perspective, un rapprochement avec l'Union IHEDN, dépositaire des fichiers des associations régionales, pourrait être envisagé.

L'enquête, dont le lancement pourrait intervenir dans le courant du premier semestre 2026, ne donnera lieu à une exploitation complète des résultats qu'au cours de l'année suivante. Elle pourrait néanmoins s'inscrire dans le cadre du cinquantième anniversaire de l'AR~16 et faire l'objet d'une restitution publique, éventuellement avec le concours de l'IRSEM, lors d'un événement organisé en 2027.

Notre ambition est réelle : l'enquête permet de célébrer notre histoire, mais aussi, dans le cadre d'une démarche participative, de réfléchir collectivement à notre avenir.

Nous vous remercions, Chère Présidente, Chères et Chers Membres du CoDir, de votre accueil, compréhension et confiance.

\begin{flushright}
Les Membres du Comité « Crises majeures »

(Présents à la réunion du 5 mars 2026)

Valérie Arlé-Faivre, Gérard Turck, Aurélien Duvignac-Rosa,\\
Jean-Jacques Cleynen, Bruno Leray, Alain Fontan
\end{flushright}

\end{tcolorbox}


%\newpage

%\printbibliography[heading=bibintoc,title={Références bibliographiques}]

%\nocite{*}

% ----------------------------------------------------------------
% Notes d'usage:
% - Le style bibliographique est dans:
%   vendor/biblatex-sciences-po-fr/style/
% - Ajouter des entrées dans les fichiers bibliographiques chargés ci-dessus:
%   vendor/biblatex-sciences-po-fr/bibliographies/{articles,livres,online,rapports,theses}.bib
% - Compilation requise:
%   latexmk -pdf main.tex
%   (latexmkrc est configuré pour XeLaTeX + Biber).
% ----------------------------------------------------------------

\end{document}
